% ---
% title: Reliability of a system
% tags: sample-space, events
% ---

\begin{problem}
	A system has 5 components, each working (1) or failed (0); the outcome is
	$(x_1,\dots,x_5)$. The system works if components 1 and 2 work, or 3 and 4
	work, or 1, 3 and 5 work. (a) What is $|\Om|$? (b) List $W$ = ``the system
	works''. (c) Let $A$ = ``components 4 and 5 failed''. What is $|A|$?
	(d) List $AW$.
	\begin{context}
		Read ``component working'' as ``issuer has not defaulted''. A
		\emph{structured product} built on issuers 1--5 keeps paying as long as at
		least one route survives, i.e.\ as long as the system works. (e) Does the
		product keep paying if any single issuer defaults? Which pairs of defaults
		stop the payments?
	\end{context}
\end{problem}

\begin{solution}
	\idea The system is three parallel routes; a route works if all its
	components work. To list $W$ without double counting, split it into
	\emph{disjoint} pieces.

	\begin{center}
		\begin{tikzpicture}[x=1cm,y=.8cm,every node/.style={font=\small},
				c/.style={draw,circle,fill=navy!10,inner sep=1.5pt}]
			\coordinate (in) at (0,0); \coordinate (out) at (5,0);
			\node[left] at (in) {in}; \node[right] at (out) {out};
			\foreach \y in {1,0,-1} {
					\draw (in) |- (.6,\y) -- (4.4,\y) -| (out);
				}
			\node[c] at (1.8,1) {1}; \node[c] at (3.2,1) {2};
			\node[c] at (1.8,0) {3}; \node[c] at (3.2,0) {4};
			\node[c] at (1.4,-1) {1}; \node[c] at (2.5,-1) {3};
			\node[c] at (3.6,-1) {5};
		\end{tikzpicture}
	\end{center}

	\step{(a)} Each $x_i\in\{0,1\}$: $|\Om|=2^5=32$.

	\step{(b)}
	\[ W=\{x_1=x_2=1\}\cup\{x_3=x_4=1\}\cup\{x_1=x_3=x_5=1\}. \]
	\begin{itemize}[nosep]
		\item Piece 1: $x_1=x_2=1$, the others free: $2^3=8$ outcomes:
		      \texttt{11000, 11001, 11010, 11011, 11100, 11101, 11110, 11111}.
		\item Piece 2: $x_3=x_4=1$ but \emph{not} $x_1=x_2=1$: $(x_1,x_2)\in
			      \{00,01,10\}$ (3 options), $x_5$ free (2 options): $6$ outcomes:
		      \texttt{00110, 00111, 01110, 01111, 10110, 10111}.
		\item Piece 3: works only through route $1,3,5$: then $x_1=x_3=x_5=1$,
		      and routes 1--2 and 3--4 are broken, so $x_2=0$ and $x_4=0$:
		      \texttt{10101}.
	\end{itemize}
	$|W|=8+6+1=15$.

	\step{(c)} $A=\{x_4=x_5=0\}$, with $x_1,x_2,x_3$ free: $|A|=2^3=8$.

	\step{(d)} If 4 and 5 failed, routes 3--4 and 1--3--5 are broken, so the
	system works only if $x_1=x_2=1$; $x_3$ is free:
	$AW=\{11000,\,11100\}$.
	\begin{answer}
		(a) $32$ \quad (b) $|W|=15$ (listed above) \quad (c) $|A|=8$ \quad
		(d) $AW=\{11000,11100\}$
	\end{answer}
	\begin{remark}
		Check: $|W^c|=32-15=17$ states in which the system fails.
	\end{remark}

	\step{(e) Finance context.} The product stops paying exactly when
	\emph{every} route contains a defaulted issuer. One default breaks at most the
	routes through that issuer: issuer 1 lies on routes 1--2 and 1--3--5 (route
	3--4 survives), issuer 3 on routes 3--4 and 1--3--5 (route 1--2 survives),
	issuers 2, 4, 5 on one route each. So any single default is survived. A pair
	of defaults must hit route $\{1,2\}$ and route $\{3,4\}$, so it is one of
	\[ \{1,3\},\quad \{1,4\},\quad \{2,3\},\quad \{2,4\}; \]
	of these, $\{1,3\}$, $\{1,4\}$ and
	$\{2,3\}$ also hit route $\{1,3,5\}$, but $\{2,4\}$ does not.
	\begin{answer}
		(e) Yes, every single default is survived; exactly the pairs $\{1,3\}$,
		$\{1,4\}$, $\{2,3\}$ stop the payments.
	\end{answer}
	\begin{remark}
		Issuers 1 and 3 appear on two routes each: they matter most for the
		product, like critical components in an engineering system.
	\end{remark}
\end{solution}
